\chapter{Od grafów plumbingowych do egzotycznych struktur gładkich: synteza}
\label{ch:tg-synteza16}

Poprzednie rozdziały wprowadzały kolejne narzędzia w różnym języku:
macierze incydencji grafów, formy przecięcia, chirurgie trójrozmaitości,
rozkłady na uchwyty i niezmienniki Seiberga--Wittena. Obecnie zestawimy
je w jeden ciąg rozumowania, pilnując, by nie wyprowadzać wniosków
gładkich z danych wyłącznie topologicznych. Rozważymy dwie konstrukcje:
wymierną kontrakcję konfiguracji $C_7$ prowadzącą do rozmaitości Parka
oraz chirurgię Fintushela--Sterna na torusie w powierzchni K3.

\section{Jaki rodzaj informacji koduje graf?}
\label{sec:tg-syn-graph}

Niech $\Gamma$ będzie skończonym, spójnym drzewem o zbiorze
wierzchołków $V(\Gamma)=\{v_1,\ldots,v_r\}$ ($r\ge1$), którego wierzchołkom
przypisano całkowite wagi $e_1,\ldots,e_r$. Zakładamy, że każda
krawędź oznacza dodatnie, poprzeczne skrzyżowanie dwóch bazowych sfer
w standardowym, zorientowanym plumbingowaniu wiązek dyskowych.
Przez $P_\Gamma$ oznaczamy wynikowe plumbingowanie, a przez
$Q_\Gamma$ macierz określoną wzorem
\[
 (Q_\Gamma)_{ij}=\begin{cases}
 e_i,&i=j,\\
 1,&i\ne j\text{ i }\{v_i,v_j\}\in E(\Gamma),\\
 0,&\text{w pozostałych przypadkach}.
 \end{cases}
\]
Macierz reprezentuje formę przecięcia na
$H_2(P_\Gamma;\mathbb Z)\cong\mathbb Z^r$ w bazie klas
obrazów przekrojów zerowych wiązek dyskowych.
Wyrazy diagonalne są liczbami Eulera wiązek
normalnych, czyli liczbami samoprzecięcia.
Poza przekątną każda krawędź daje jeden punkt
przecięcia ze znakiem $+1$, a brak krawędzi
oznacza rozłączność odpowiednich sfer.
To stwierdzenie zależy od określonego sposobu plumbingowania;
nie należy utożsamiać dowolnego grafu ważonego z jednoznacznie
wyznaczoną czterorozmaitością bez podania danych konstrukcji.

\begin{twierdzenie}[Homologia brzegu niezdegenerowanego plumbingowania]
\label{thm:tg-syn-boundary}
Niech $\Gamma$ będzie drzewem z całkowitymi wagami, a
$P_\Gamma$ standardowym plumbingowaniem zorientowanych wiązek
dyskowych nad $S^2$. Niech $Y=\partial P_\Gamma$ z orientacją
brzegową i załóżmy, że $\det Q_\Gamma\ne0$. Wówczas
$H_1(Y;\mathbb Z)\cong\operatorname{coker}Q_\Gamma$,
$|H_1(Y;\mathbb Z)|=|\det Q_\Gamma|$ i
$H_1(Y;\mathbb Q)=0$. Jeśli $Q_\Gamma$ jest ujemnie określona,
to $\sigma(P_\Gamma)=-r$.
\end{twierdzenie}

\begin{proof}
Drzewo plumbingowe daje rozkład na uchwyty z jednym uchwytem indeksu
zero i $r$ uchwytami indeksu dwa, bez uchwytów indeksu jeden.
Stąd $P_\Gamma$ jest po prostu spójna,
$H_2(P_\Gamma;\mathbb Z)\cong\mathbb Z^r$ i
$H_1(P_\Gamma;\mathbb Z)=0$. W długim ciągu dokładnym pary
$(P_\Gamma,Y)$ otrzymujemy
\[
 H_2(P_\Gamma;\mathbb Z)\xrightarrow{j}
 H_2(P_\Gamma,Y;\mathbb Z)\longrightarrow
 H_1(Y;\mathbb Z)\longrightarrow0.
\]
Z dualności Poincar\'ego--Lefschetza i uniwersalnego twierdzenia
o współczynnikach wynika
$H_2(P_\Gamma,Y;\mathbb Z)\cong H^2(P_\Gamma;\mathbb Z)
\cong\operatorname{Hom}(H_2(P_\Gamma;\mathbb Z),\mathbb Z)$.
W tych identyfikacjach $j$ jest homomorfizmem wyznaczonym przez
formę przecięcia, a więc ma macierz $Q_\Gamma$. Jego kokernel
jest zatem grupą $H_1(Y;\mathbb Z)$. Niezerowy wyznacznik
implikuje, że kokernel jest skończony. Istotnie,
postać normalna Smitha macierzy $Q_\Gamma$ ma
niezerowe dodatnie wyrazy diagonalne $d_1,\ldots,d_r$;
kokernel jest sumą prostą grup $\mathbb Z/d_i\mathbb Z$,
a iloczyn $d_1\cdots d_r$ wynosi $|\det Q_\Gamma|$,
ponieważ operacje unimodularne zmieniają wyznacznik
co najwyżej o znak. Po przejściu do $\mathbb Q$ mapa staje się
izomorfizmem, więc $H_1(Y;\mathbb Q)=0$.
Brzeg jest spójny (ma opis jako chirurgia na splocie
w $S^3$), zamknięty i zorientowany; dualność
Poincarégo daje również $H_2(Y;\mathbb Q)=0$.
Tak więc $Y$ jest wymierną sferą homologiczną.
Jeżeli forma jest ujemnie określona, wszystkie jej $r$
wartości własnych są ujemne, stąd $\sigma(P_\Gamma)=-r$.
\end{proof}

\begin{uwaga}[Forma splecenia na brzegu]
W hipotezach twierdzenia i przy orientacji $Y=\partial P_\Gamma$
forma splecenia na $H_1(Y;\mathbb Z)$ jest reprezentowana
przez $-Q_\Gamma^{-1}$ modulo $\mathbb Z$. Dokładniej,
w prezentacji $\mathbb Z^r/Q_\Gamma\mathbb Z^r$ mamy
\[
 \lambda_Y([x],[y])=-x^TQ_\Gamma^{-1}y+\mathbb Z.
\]
Identyfikację z topologiczną formą splecenia
przyjmujemy z rachunku chirurgicznego; algebraicznie
wzór jest poprawnie określony, ponieważ zastąpienie
$x$ przez $x+Q_\Gamma z$ zmienia wartość o
$-z^Ty\in\mathbb Z$, a analogicznie dla drugiego
argumentu. Forma jest symetryczna i niezdegenerowana:
jeżeli $x^TQ_\Gamma^{-1}y$ jest całkowite dla
każdego $y\in\mathbb Z^r$, to $Q_\Gamma^{-1}x$
ma całkowite współrzędne, czyli $[x]=0$.
Zmiana orientacji $Y$ zmienia znak tej formy.
Jest to ta sama konwencja, której używamy w stwierdzeniu
\ref{prop:tg-rb-linkingform} i w lemacie
\ref{lem:tg-park-metabolizer}. W szczególności dla
pierwszego meridianu konfiguracji $C_7$ otrzymujemy
$\lambda_{Y_7}([\mu_0],[\mu_0])=6/49\pmod{\mathbb Z}$,
a nie $-6/49$. Własności algebraiczne takie jak
izotropowość podgrupy $7\mathbb Z/49\mathbb Z$
pozostają niezmienione po zmianie orientacji.
Sama grupa $H_1(Y)$, a nawet
forma splecenia, nie wyznaczają ogólnie typu dyfeomorfizmu $Y$.
Na przykład rozpoznanie konkretnej przestrzeni soczewkowej
wymaga dodatkowo informacji o sposobie sklejenia torusów pełnych.
\end{uwaga}

\begin{przyklad}[Łańcuch $C_p$ i ułamek ciągły]
\label{ex:tg-syn-cp}
Niech $p\ge2$ i niech $C_p$ będzie łańcuchem długości $p-1$
o wagach $(-(p+2),-2,\ldots,-2)$. Dla wyznacznika
$D_k$ lewego głównego minora rozmiaru $k$ macierzy przecięcia
mamy $D_0=1$, $D_1=-(p+2)$ i
$D_k=-2D_{k-1}-D_{k-2}$ dla $2\le k\le p-1$.
Rekurencję otrzymujemy przez rozwinięcie wyznacznika
wzdłuż ostatniego wiersza macierzy trójdiagonalnej:
wkład ostatniego wyrazu diagonalnego wynosi
$-2D_{k-1}$, a wkład dwóch jednostek
poza przekątną wynosi $-D_{k-2}$.
Indukcyjnie otrzymujemy
$D_k=(-1)^k(1+k(p+1))$: wzór zachodzi dla $k=0,1$,
a podstawienie do rekurencji daje
\[
 -2(-1)^{k-1}(1+(k-1)(p+1))
 -(-1)^{k-2}(1+(k-2)(p+1))
 =(-1)^k(1+k(p+1)).
\]
Stąd $\det Q_{C_p}=D_{p-1}=(-1)^{p-1}p^2$.
Ponieważ $(-1)^kD_k>0$ dla każdego $k$,
wszystkie wiodące główne minory macierzy $-Q_{C_p}$
są dodatnie. Kryterium Sylvestera daje dodatnią
określoność $-Q_{C_p}$, a więc ujemną
określoność $Q_{C_p}$.
Z kolei ujemny ułamek łańcuchowy opisujący brzeg ma postać
\[
[p+2,\underbrace{2,\ldots,2}_{p-2\text{ razy}}]^-
=p+2-\cfrac{1}{2-\cfrac{1}{\ddots-\cfrac12}}
=\frac{p^2}{p-1}.
\]
Aby uzasadnić ostatnią równość, oznaczmy przez $f_k$
ujemny ułamek łańcuchowy złożony z $k$ dwójek.
Mamy $f_1=2$ i $f_{k+1}=2-1/f_k$.
Jeżeli $f_k=(k+1)/k$, to
$f_{k+1}=2-k/(k+1)=(k+2)/(k+1)$.
Stąd indukcyjnie $f_k=(k+1)/k$.
Dla $p\ge3$, po podstawieniu $k=p-2$, dostajemy
$p+2-(p-2)/(p-1)=p^2/(p-1)$.
Dla $p=2$ lista dwójek jest pusta i ułamek wynosi $4$,
co również zgadza się z wzorem.
Ten drugi rachunek, a nie sama wartość wyznacznika,
pozwala zidentyfikować brzeg jako przestrzeń
soczewkową $L(p^2,1-p)$ przy konwencji orientacji
z twierdzenia \ref{thm:tg-rb-lens}.
\end{przyklad}

\begin{lemat}[Grupa dyskryminantowa i metabolizer łańcucha $C_p$]
\label{lem:tg-syn-cp-metabolizer}
Niech $p\ge2$ będzie liczbą całkowitą,
$C_p$ standardowym łańcuchem o wagach
$(-(p+2),-2,\ldots,-2)$ i $Y_p=\partial C_p$.
Niech $\mu_0$ oznacza meridian pierwszej składowej
ramowanego splotu, o ramowaniu $-(p+2)$.
Wtedy
\[
H_1(Y_p;\mathbb Z)=\langle[\mu_0]\rangle
\cong\mathbb Z/p^2\mathbb Z.
\]
Przy konwencji formy splecenia
$\lambda_{Y_p}=-Q_{C_p}^{-1}\pmod{\mathbb Z}$
mamy
\[
\lambda_{Y_p}([\mu_0],[\mu_0])
=\frac{p-1}{p^2}\pmod{\mathbb Z}.
\]
Podgrupa $M_p:=p\langle[\mu_0]\rangle$
jest rzędu $p$ i spełnia $M_p=M_p^\perp$,
gdzie dopełnienie ortogonalne jest liczone
względem formy splecenia w grupie $H_1(Y_p;\mathbb Z)$.
Jeśli dodatkowo istnieje wymierna kula homologii
$B_p$ z $\partial B_p\cong Y_p$,
$H_1(B_p;\mathbb Z)\cong\mathbb Z/p\mathbb Z$
oraz surjektywnym odwzorowaniem
$H_1(Y_p;\mathbb Z)\to H_1(B_p;\mathbb Z)$,
to $M_p$ jest jądrem tego odwzorowania.
\end{lemat}

\begin{proof}
Niech $n=p-1$. W prezentacji grupy
$H_1(Y_p;\mathbb Z)=\operatorname{coker}Q_{C_p}$
generatory oznaczmy $x_0,\ldots,x_{n-1}$.
Gdy $n=1$, jedyna relacja ma postać
$-4x_0=0$, a teza o cykliczności jest oczywista.
Dla $n\ge2$ ostatnia relacja daje
$x_{n-2}=2x_{n-1}$, a kolejne relacje
$ x_{i-1}-2x_i+x_{i+1}=0$
pozwalają indukcyjnie wyprowadzić
$x_i=(n-i)x_{n-1}$ dla $0\le i\le n-1$.
Pierwsza relacja przyjmuje wówczas postać
\[
-(p+2)(p-1)x_{n-1}+(p-2)x_{n-1}
=-p^2x_{n-1}=0.
\]
Wszystkie generatory są wielokrotnościami
$x_{n-1}$, zaś rząd grupy jest równy
$|\det Q_{C_p}|=p^2$, więc
$H_1(Y_p;\mathbb Z)\cong\mathbb Z/p^2\mathbb Z$.
Ponieważ $x_0=(p-1)x_{n-1}$ i
$\gcd(p-1,p^2)=1$, także $x_0=[\mu_0]$
jest generatorem (dla $p=2$ to ten sam generator).

Główny minor powstały po skreśleniu pierwszego
wiersza i kolumny $Q_{C_p}$ jest macierzą
łańcucha $n-1$ dwójek ujemnych;
jego wyznacznik wynosi
$(-1)^{n-1}n$. Dokładniej, wyznaczniki $F_k$
łańcuchów $k$ ujemnych dwójek spełniają
$F_0=1$, $F_1=-2$, $F_k=-2F_{k-1}-F_{k-2}$.
Podstawienie $F_k=(-1)^k(k+1)$ do rekurencji
i sprawdzenie warunków początkowych dowodzą
wzoru również dla pustej macierzy, gdy $n=1$.
Wzór na element odwrotności
oraz $\det Q_{C_p}=(-1)^np^2$
dają $(Q_{C_p}^{-1})_{00}=-(p-1)/p^2$,
stąd podany wzór dla $\lambda$.
Można także sprawdzić całą pierwszą kolumnę:
dla $v=(n,n-1,\ldots,1)^T$ zachodzi
$Q_{C_p}v=-p^2 e_0$, gdzie $e_0$ jest pierwszym
wektorem standardowej bazy $\mathbb Z^n$.
Ostatnia współrzędna iloczynu znika dla $n\ge2$,
współrzędne wewnętrzne znikają na mocy
stałości różnic kolejnych wyrazów, a pierwsza
wynosi $-(p+2)(p-1)+(p-2)=-p^2$.
Dla $n=1$ sprawdzamy bezpośrednio $(-4)(1)=-4$.
Wynika stąd $Q_{C_p}^{-1}e_0=-v/p^2$.

Element $p[\mu_0]$ ma rząd $p$ i
$\lambda(p[\mu_0],p[\mu_0])=p-1=0$
w $\mathbb Q/\mathbb Z$.
Ponadto $x[\mu_0]$ jest ortogonalny
do $M_p$ dokładnie wtedy, gdy
$(p-1)x/p$ jest liczbą całkowitą.
Ponieważ $\gcd(p-1,p)=1$,
jest to równoważne $p\mid x$.
Stąd $M_p^\perp=M_p$.

Przy dodatkowych hipotezach o $B_p$
jądro wspomnianej surjekcji ma rząd $p$.
W cyklicznej grupie rzędu $p^2$
istnieje dokładnie jedna podgrupa rzędu $p$,
czyli $M_p$.
W szczególności standardowa kula Cassona--Harera
spełnia te hipotezy na mocy twierdzenia
\ref{thm:tg-rb-casson-harer}: surjekcja grup
fundamentalnych indukuje surjekcję ich abelianizacji,
czyli pierwszych grup homologii. Istnienie metabolizera
jest warunkiem algebraicznym, a nie samodzielnym
dowodem istnienia gładkiego wymiernego wypełnienia.
\end{proof}

\section{Dwie operacje wymiany: co rzeczywiście zachowują?}
\label{sec:tg-syn-exchange}

W każdej z dwóch omawianych operacji czterorozmaitość
jest sumą dwóch zwartych części o wspólnym trójwymiarowym
brzegu. Zastępujemy jedną część inną, nie zmieniając dopełnienia.
Dane potrzebne do obliczeń można zorganizować następująco.

\begin{center}
\begin{tikzpicture}[>=Stealth,line cap=round,line join=round,
    every node/.style={align=center}]
\node[draw=tgCoral,fill=tgCoral!10,rounded corners=4pt,
    text width=4.3cm,inner sep=6pt] (p) at (-3.25,1.2)
    {\small \textbf{wymierna kontrakcja}\\
     \footnotesize $C_p\rightsquigarrow B_p$\\
     \footnotesize wspólny brzeg $L(p^2,1-p)$};
\node[draw=tgTeal,fill=tgTeal!10,rounded corners=4pt,
    text width=4.3cm,inner sep=6pt] (k) at (3.25,1.2)
    {\small \textbf{chirurgia wzdłuż węzła}\\
     \footnotesize $T^2\times D^2\rightsquigarrow S^1\times E_K$\\
     \footnotesize wspólny brzeg $T^3$};
\node[draw=tgPurple,fill=tgPurple!9,rounded corners=4pt,
    text width=9cm,inner sep=6pt] (d) at (0,-1.05)
    {\small zachowane dopełnienie $A$ i identyfikacja brzegu\\
     \footnotesize Mayer--Vietoris, van Kampen oraz formy przecięcia};
\draw[-{Stealth[length=2.7mm]},tgPurple,line width=1.5pt] (p.south)--(-2.2,-.64);
\draw[-{Stealth[length=2.7mm]},tgPurple,line width=1.5pt] (k.south)--(2.2,-.64);
\end{tikzpicture}
\end{center}

\begin{twierdzenie}[Bilans wymiernej kontrakcji]
\label{thm:tg-syn-rational}
Niech $X$ będzie gładką, zamkniętą, spójną i zorientowaną
czterorozmaitością zawierającą standardową konfigurację
plumbingową $C_p$ o $p-1$ wierzchołkach. Niech $A=X\setminus\operatorname{int}C_p$.
Niech $B_p$ będzie zorientowaną wymierną kulą
homologii, której brzeg jest dyfeomorficzny
z $\partial C_p$ z zachowaniem orientacji.
Dołączamy ją do $A$ przez dyfeomorfizm
$\partial B_p\to\partial A$ odwracający orientację
brzegową (orientacja $\partial A$ jest przeciwna
do orientacji $\partial C_p$). Oznaczmy wynik
przez $X_{(p)}$. Obie rozmaitości wynikowe są
zamknięte i spójne: brzeg wycinanej części jest
spójny, a każda składowa dopełnienia musi go spotykać
ze względu na spójność $X$. Dołączenie $B_p$
nie zmienia tej własności.
Wówczas
\[
 \chi(X_{(p)})=\chi(X)-(p-1),\qquad
 \sigma(X_{(p)})=\sigma(X)+(p-1).
\]
Ponadto bez dodatkowych założeń na $b_1$ zachodzą
równości $b_2^+(X_{(p)})=b_2^+(X)$ oraz
$b_2^-(X_{(p)})=b_2^-(X)-(p-1)$,
jak wynika z następującego wniosku.
\end{twierdzenie}

\begin{proof}
Plumbingowanie $C_p$ ma rozkład na uchwyty z jednym uchwytem
indeksu zero i $p-1$ uchwytami indeksu dwa,
zatem $\chi(C_p)=p$. Wymierna kula homologii ma
homologię wymierną punktu, więc $\chi(B_p)=1$.
Przy sklejeniu wzdłuż zamkniętej trójrozmaitości $Y$
zachodzi $\chi(X)=\chi(A)+\chi(C_p)-\chi(Y)$ i analogicznie
dla $X_{(p)}$, przy czym $\chi(Y)=0$.
Po odjęciu obu wzorów dostajemy
$\chi(X_{(p)})-\chi(X)=1-p$.
Forma przecięcia $C_p$ jest ujemnie określona rzędu $p-1$,
więc $\sigma(C_p)=-(p-1)$. Forma przecięcia $B_p$
ma zerowy wymiar nad $\mathbb Q$, dlatego $\sigma(B_p)=0$.
Z addytywności sygnatury Nowikowa przy sklejaniu otrzymujemy
drugi wzór.

Aby przejść od tych równości do liczb Betti'ego,
wystarcza zachowanie $b_1$, ponieważ dla dowolnej
zamkniętej, spójnej, zorientowanej czterorozmaitości
$\chi=2-2b_1+b_2$.
Zachowanie $b_1$ i dokładniejszy rozkład formy
przecięcia uzasadnimy w następującym wniosku.
\end{proof}

\begin{wniosek}[Wymierny bilans bez założenia o pierwszej homologii]
\label{cor:tg-syn-rational-homological}
W hipotezach Twierdzenia~\ref{thm:tg-syn-rational},
zachodzą w szczególności
\[
 b_1(X_{(p)})=b_1(X),\qquad
 b_2^+(X_{(p)})=b_2^+(X),\qquad
 b_2^-(X_{(p)})=b_2^-(X)-(p-1).
\]
Ponadto jako formy symetryczne nad $\mathbb Q$
zachodzą izometrie
\[
 Q_X\cong_{\mathbb Q}Q_{X_{(p)}}\perp Q_{C_p},
 \qquad \sigma(Q_{C_p})=-(p-1).
\]
Nie oznacza to istnienia analogicznego rozkładu
\emph{całkowitej} formy unimodularnej.
W istocie bezpośrednia suma ortogonalna całkowitych
form unimodularnej rozmaitości $X_{(p)}$ i $C_p$
miałaby wyznacznik o module $p^2>1$,
podczas gdy forma $X$ na homologii modulo torsję
ma wyznacznik o module jeden. Dla każdej zamkniętej
zorientowanej czterorozmaitości wynika to z dualności
Poincarégo i uniwersalnego twierdzenia o współczynnikach:
surjekcja $H^2\to\operatorname{Hom}(H_2,\mathbb Z)$
ma skończone jądro $\operatorname{Ext}^1_{\mathbb Z}(H_1,\mathbb Z)$,
równe podgrupie torsyjnej $H^2$.
Po złożeniu z dualnością i przejściu modulo torsję
otrzymujemy izomorfizm wyznaczony przez formę przecięcia. Na poziomie
całkowitym występują więc dodatkowe dane
sklejenia przez torsję brzegu.
\end{wniosek}

\begin{proof}
Oznaczmy $A=X\setminus\operatorname{int}C_p$
oraz $Y=\partial C_p$. Ponieważ $Q_{C_p}$
ma niezerowy wyznacznik, $Y$ jest
wymierną sferą homologiczną, a zatem
$H_1(Y;\mathbb Q)=H_2(Y;\mathbb Q)=0$.
Wymierna kula homologii $B_p$
spełnia $H_1(B_p;\mathbb Q)
=H_2(B_p;\mathbb Q)=0$.
Z długich ciągów Mayera--Vietorisa
dla $X=A\cup_YC_p$ i
$X_{(p)}=A\cup_YB_p$
otrzymujemy izomorfizmy
\[
 H_2(X;\mathbb Q)\cong
 H_2(A;\mathbb Q)\oplus H_2(C_p;\mathbb Q),
 \qquad H_2(X_{(p)};\mathbb Q)
 \cong H_2(A;\mathbb Q).
\]
Ponieważ $Y,C_p,B_p$ są spójne, strzałki
z $H_0(Y;\mathbb Q)$ do sum grup $H_0$
sklejanych części są injektywne (już rzut na
$H_0(C_p;\mathbb Q)$ lub $H_0(B_p;\mathbb Q)$
jest izomorfizmem). W stopniu jeden mamy
ponadto $H_1(Y;\mathbb Q)=0$ oraz
$H_1(C_p;\mathbb Q)=H_1(B_p;\mathbb Q)=0$.
Stąd ciągi Mayera--Vietorisa dają $H_1(X;\mathbb Q)\cong
H_1(A;\mathbb Q)
\cong H_1(X_{(p)};\mathbb Q)$.

Kołnierze brzegu pozwalają przesunąć cykle
z obu części do ich wnętrz bez zmiany klas.
Ich przecięcia wzajemne są zerowe, ponieważ
wnętrza są rozłączne. Wewnętrzne parowania
w każdej części obliczamy za pomocą mapy
$H_2(W;\mathbb Q)\to H_2(W,\partial W;\mathbb Q)$
oraz dualności Poincarégo--Lefschetza.
Dla $W=A,C_p$ mapa ta jest izomorfizmem:
w ciągu pary po obu jej stronach występują
zerowe grupy $H_2(Y;\mathbb Q)$ i $H_1(Y;\mathbb Q)$.
Parowanie części $A$ nie zależy zatem od wyboru
wypełnienia, a parowanie części $C_p$ to $Q_{C_p}$.
Rozkład stopnia dwa jest więc
ortogonalny względem formy
przecięcia i daje wskazaną
izometrię nad $\mathbb Q$.
Ponieważ $Q_{C_p}$ jest ujemnie
określona rzędu $p-1$,
dodaje dokładnie $p-1$ ujemnych
kierunków i żadnego dodatniego.
To dowodzi wszystkich równości.
\end{proof}

\begin{stwierdzenie}[Bilans chirurgii wzdłuż węzła]
\label{prop:tg-syn-knot}
Niech $X$ będzie spójną, zamkniętą, gładką i zorientowaną
czterorozmaitością z gładko zanurzonym torusem $T$
o trywialnej zorientowanej wiązce normalnej.
Niech $X_K$ będzie wynikiem standardowego sklejenia
z rozdziału \ref{ch:tg-knot-surgery}. Wówczas
$\chi(X_K)=\chi(X)$ i $\sigma(X_K)=\sigma(X)$.
Przy tym sklejeniu całkowite grupy homologii $X_K$ i $X$
są izomorficzne, a całkowite formy przecięcia
na drugiej homologii modulo torsję są izometryczne.
\end{stwierdzenie}

\begin{proof}
Dla $P_K=S^1\times E_K$ z poprzedniego rozdziału mamy
$H_i(P_K;\mathbb Z)\cong H_i(T^2\times D^2;\mathbb Z)$
we wszystkich stopniach. Co więcej, ustalone sklejenie
identyfikuje odwzorowania z $H_i(T^3)$ do homologii tych dwóch
części. Dla dwóch rozkładów z dopełnieniem $A$ ciągi
Mayera--Vietorisa dają więc w każdym stopniu krótkie ciągi
\[
 0\longrightarrow\operatorname{coker} j_i
 \longrightarrow H_i(X_K;\mathbb Z)
 \longrightarrow\ker j_{i-1}\longrightarrow0,
\]
gdzie $j_i$ jest parą map indukowanych przez
włączenia brzegu do obu części (z przeciwnymi
znakami), o wartościach w
$H_i(A;\mathbb Z)\oplus H_i(P_K;\mathbb Z)$.
Analogiczne krótkie ciągi występują dla $X$ i $T^2\times D^2$.
Ponieważ $H_{i-1}(T^3;\mathbb Z)$ jest grupą wolną
(przy $i=0$ przyjmujemy $H_{-1}=0$), także
$\ker j_{i-1}$ jest wolna. Każdy taki ciąg rozkłada się
jako ciąg grup abelowych. Identyfikacja jąder i kokerneli
zapewnia zatem izomorfizmy całkowitych grup homologii.

Obie części dołączane oraz wspólny brzeg mają
zerową charakterystykę Eulera, więc jej addytywność
daje $\chi(X_K)=\chi(X)$. Ich formy przecięcia
mają zerową sygnaturę (generatory drugich grup
homologii pochodzą z brzegu), zatem addytywność
sygnatury Nowikowa daje $\sigma(X_K)=\sigma(X)$.
Samo istnienie izomorfizmów grup homologii
nie ustanawia jednak izometrii form przecięcia.
W uzupełnieniu dowodu po wniosku
\ref{cor:tg-ks-invariants} skonstruowaliśmy
odwzorowanie stopnia jeden $p:X_K\to X$,
równe identyczności na $A$ i indukujące
izomorfizm homologii całkowitej.
Naturalność iloczynu w kohomologii i dualność
Poincarégo dają izometrię $p_*$ oraz
$p_*[T_K]=[T]$. To jest brakujący krok
przy przejściu od homologii do formy przecięcia.
\end{proof}

\section{Grupa fundamentalna jako kontrola po sklejeniu}
\label{sec:tg-syn-pi1}

Gdy $X=A\cup_Y W$, twierdzenie van Kampena daje prezentację
$\pi_1(X)$ jako wypchnięcia grup $\pi_1(A)$ i
$\pi_1(W)$ względem dwóch homomorfizmów
z $\pi_1(Y)$, przy spójnych częściach i brzegu.
Nie zakładamy injektywności tych homomorfizmów:
w ilorazie wolnego iloczynu narzucamy równości
obrazów każdego elementu $\pi_1(Y)$. Poniższe dwa zastosowania
wyjaśniają, dlaczego samo obliczenie $H_1(Y)$ nie wystarcza.

\begin{stwierdzenie}[Prosta spójność po chirurgii wzdłuż węzła]
\label{prop:tg-syn-pi1}
Niech $X$ będzie spójną, zamkniętą, gładką czterorozmaitością,
$T\subset X$ torusem o trywialnej wiązce normalnej i niech
$A=X\setminus\operatorname{int}\nu T$ będzie spójna i po prostu
spójna. Wówczas $\pi_1(X_K)=1$ dla dowolnego węzła
$K\subset S^3$ i standardowego sklejenia.
\end{stwierdzenie}

\begin{proof}
Mamy $\pi_1(S^1\times E_K)\cong\mathbb Z\times\pi_1(E_K)$.
Przy $\pi_1(A)=1$ twierdzenie van Kampena daje
\[
 \pi_1(X_K)\cong
 \frac{\mathbb Z\times\pi_1(E_K)}
 {\langle\!\langle s,\mu_K,\lambda_K\rangle\!\rangle}.
\]
Meridian $\mu_K$ normalnie generuje $\pi_1(E_K)$: istotnie
meridionalne wypełnienie Dehna odtwarza $S^3$, więc
$\pi_1(E_K)/\langle\!\langle\mu_K\rangle\!\rangle=1$.
Ponadto relacja $s=1$ usuwa pozostały czynnik $\mathbb Z$.
Zatem grupa wynikowa jest trywialna.
\end{proof}

Dla wymiernej kontrakcji argument jest inny:
$\pi_1(B_p)$ i $\pi_1(\partial C_p)$ mogą być nietrywialne.
Aby uzyskać prostą spójność, trzeba sprawdzić obrazy ich
generatorów w dopełnieniu konfiguracji. W konstrukcji Parka
niezbędny fakt geometryczny polega na wskazaniu
generatora cyklicznej grupy $\pi_1(\partial C_7)$,
który ogranicza dysk w dopełnieniu.
W dowodzie twierdzenia \ref{thm:tg-park-pi1}
jest to meridian $\mu_3=3\mu_5$ (w zapisie
addytywnym grupy cyklicznej rzędu $49$).
Ponieważ $\gcd(3,49)=1$, zanika obraz całej grupy
brzegu. Ponieważ $R$ i $C_7$ są po prostu spójne,
van Kampen daje $\pi_1(R)\cong
\pi_1(A)/\langle\!\langle i_*(\pi_1(Y_7))\rangle\!\rangle$.
Obraz $i_*(\pi_1(Y_7))$ jest trywialny, więc
$\pi_1(A)=1$. Następnie surjekcja
$\pi_1(\partial B_7)\to\pi_1(B_7)$ pozwala
wywnioskować $\pi_1(X_7)=1$. Bez tego
kroku sam rachunek $\det Q_{C_7}=49$ nie wyznacza
$\pi_1(X_7)$.

\section{Formy przecięcia i klasyfikacja topologiczna}
\label{sec:tg-syn-freedman}

\begin{twierdzenie}[Obliczenie typu topologicznego rozmaitości Parka]
\label{thm:tg-syn-park-homeo}
Niech $X_7$ będzie gładką, zamkniętą, spójną,
zorientowaną i po prostu spójną czterorozmaitością
uzyskaną z $R=E(1)\#4\overline{\mathbb{CP}}^{\,2}$
przez standardową wymierną kontrakcję konfiguracji $C_7$.
Wtedy jej forma przecięcia jest nieparzysta, unimodularna
i nieokreślona, o sygnaturze $-6$ i randze $8$.
W konsekwencji, przy wykorzystaniu twierdzenia Freedmana,
$X_7$ jest homeomorficzna z
$\mathbb{CP}^2\#7\overline{\mathbb{CP}}^{\,2}$.
\end{twierdzenie}

\begin{proof}
Ponieważ $E(1)\cong\mathbb{CP}^2\#9\overline{\mathbb{CP}}^{\,2}$,
mamy $\chi(E(1))=12$ oraz $\sigma(E(1))=-8$.
Rozdmuchanie czterech punktów zwiększa charakterystykę Eulera
o $4$, a zmniejsza sygnaturę o $4$; stąd
$\chi(R)=16$, $\sigma(R)=-12$.
Konfiguracja $C_7$ ma sześć sfer, więc
z twierdzenia \ref{thm:tg-syn-rational} dostajemy
$\chi(X_7)=16-6=10$, $\sigma(X_7)=-12+6=-6$.
Z prostej spójności wynika $b_1(X_7)=0$ i
$\chi(X_7)=2+b_2(X_7)$, zatem $b_2(X_7)=8$.
Stąd $b_2^+=1$ i $b_2^-=7$. Forma przecięcia
na $H_2(X_7;\mathbb Z)$ jest unimodularna
z dualności Poincar\'ego, gdyż $H_1(X_7;\mathbb Z)=0$
i nie występuje torsja w $H_2$:
uniwersalne twierdzenie o współczynnikach daje
$H^2\cong\operatorname{Hom}(H_2,\mathbb Z)$,
a dualność Poincarégo identyfikuje $H_2$ z $H^2$.
Złożenie tych izomorfizmów jest dokładnie
homomorfizmem $x\mapsto Q_{X_7}(x,-)$.
Ponieważ homologia zwartej rozmaitości jest
skończenie generowana, grupa homomorfizmów
$\operatorname{Hom}(H_2,\mathbb Z)$ jest wolna;
stąd wolność $H_2$ oraz unimodularność macierzy
tego homomorfizmu.
Forma nie może być parzysta: korzystamy z klasycznego
twierdzenia o sygnaturze parzystych unimodularnych
form całkowitych (jest podzielna przez $8$),
a tutaj sygnatura wynosi $-6$. Alternatywny argument
oparty na twierdzeniu Rokhlina podano w lemacie
\ref{lem:tg-park-odd}. Z klasyfikacji nieparzystych
nieokreślonych form unimodularnych wynika więc
$Q_{X_7}\cong\langle1\rangle\oplus7\langle-1\rangle$.
Obie rozmaitości są gładkie, więc oba niezmienniki
Kirby'ego--Siebenmanna wynoszą zero. W tym miejscu stosujemy
głębokie twierdzenie Freedmana o klasyfikacji po prostu spójnych
zamkniętych topologicznych czterorozmaitości i otrzymujemy
też zgodność typu homeomorfizmu z rozmaitością standardową.
\end{proof}

\begin{uwaga}[Dlaczego nie wystarcza sam wzór na $\chi$ i $\sigma$?]
W dowodzie wykorzystano dodatkowo prostą spójność, dualność
Poincar\'ego, parzystość formy i klasyfikację unimodularnych
form całkowitych. Z dwóch liczb $\chi$ i $\sigma$ bez tych
hipotez nie można wyprowadzić homeomorfizmu.
\end{uwaga}

\begin{przyklad}[Topologiczny model chirurgii na K3]
Powierzchnia K3 ma $\chi=24$, $\sigma=-16$,
$b_2^+=3$, $b_2^-=19$ oraz parzystą unimodularną
formę przecięcia $2(-E_8)\oplus3H$, gdzie
$H=\begin{psmallmatrix}0&1\\1&0\end{psmallmatrix}$.
Ze stwierdzeń \ref{prop:tg-syn-knot}
i \ref{prop:tg-syn-pi1}, zastosowanych do włókna
o jednospójnym dopełnieniu, wynika, że wszystkie
$E(2)_K$ mają tę samą formę przecięcia i są
po prostu spójne. Istnienie takiego włókna
uzasadniono w sekcji \ref{sec:tg-ks-examples}.
Twierdzenie Freedmana daje ich homeomorfizm z K3.
Sama równość $\chi=24$ i $\sigma=-16$ nie gwarantowałaby
parzystości formy ani takiego wniosku.
\end{przyklad}

\section{Dwa mechanizmy wykrywania różnic gładkich}
\label{sec:tg-syn-sw}

W tych dwóch konstrukcjach powyższe dane topologiczne
nie wystarczają do rozróżnienia struktur gładkich.
Potrzebujemy dodatkowych niezmienników.
Poniżej porównamy dwa sposoby użycia niezmienników
Seiberga--Wittena.

\begin{stwierdzenie}[Dwa reżimy niezmiennika SW]
\label{prop:tg-syn-regimes}
Niech $X$ będzie gładką, zamkniętą, zorientowaną
czterorozmaitością z $b_1(X)=0$.
Ustalmy orientację homologiczną i strukturę
$\mathrm{Spin}^c$ o wymiarze oczekiwanym $d=0$.
Jeżeli $b_2^+(X)>1$, zliczanie rozwiązań równań
Seiberga--Wittena dla generycznych danych wyznacza
niezmienniki niezależne od wyboru komory. Jeśli natomiast
$b_2^+(X)=1$, w ogólności trzeba wskazać komorę
danych metrycznych i perturbacji; wartości w dwóch
komorach mogą się różnić. W szczególności bezkomorowa
formuła Fintushela--Sterna nie może być automatycznie
zastosowana do rozmaitości Parka $X_7$.
\end{stwierdzenie}

\begin{proof}[Uzasadnienie zakresu]
Na mocy stwierdzenia \ref{prop:tg-sw-wall}
dla ustalonej metryki ściana redukowalnych jest
poziomicą surjektywnego rzutu afinicznego na przestrzeń
harmonicznych form samodualnych wymiaru $b_2^+$.
Przy zmianie metryki przestrzenie te tworzą
wiązkę wektorową tej samej rangi. Przy $b_2^+>1$
generyczną jednoparametrową rodzinę danych można
więc wybrać tak, by unikała ściany, stosując
transwersalność względem jej przekroju zerowego. Nie ma wtedy zmiany komory
podczas porównania dwóch generycznych danych.
Przy $b_2^+=1$ zbiór takich danych ma kowymiar jeden:
może rozdzielać przestrzeń danych, a przekraczanie ściany może
zmienić zorientowane zliczenie rozwiązań.
Dokładne uzasadnienie niezmienniczości wewnątrz
komory podano w stwierdzeniu \ref{prop:tg-sw12-count}.
Lokalny mechanizm zmiany komory, wraz z jego
dodatkowymi hipotezami, przedstawiono w twierdzeniu
\ref{thm:tg-sw12-wall-model}. Nie zakładamy, że
każde przecięcie ściany jest automatycznie regularne.
\end{proof}

\begin{przyklad}[Park: jeden niezerowy niezmiennik wystarcza]
Dla $X_7$ mamy $b_2^+=1$. W skrypcie obliczono
$K_{X_7}^2=2$ i udowodniono, przy wykorzystaniu
geometrycznego wyboru formy symplektycznej, że
$K_{X_7}\cdot[\omega_7]>0$. Dla kanonicznej struktury $\mathrm{Spin}^c$
klasa wyznacznikowa wynosi $c=-K_{X_7}$.
Ponieważ $c^2>0$, komora małych perturbacji
nie zależy od metryki (lemat
\ref{lem:tg-park14-no-wall}). Nierówność
$c\cdot[\omega_7]<0$ identyfikuje ją z komorą
Taubesa--Li--Liu, w której
$SW_{X_7}^{\circ}(c)=\pm1$.
Tożsamość obu komór jest zasadniczym krokiem
twierdzenia \ref{thm:tg-park14-exotic}. Argument z metryką
o dodatniej krzywiźnie skalarnej wyklucza dyfeomorfizm
ze standardową $\mathbb{CP}^2\#7\overline{\mathbb{CP}}^{\,2}$.
Nie jest to konsekwencja samej formy przecięcia, lecz
głębokiej teorii gładkiej i poprawnego doboru komory.
\end{przyklad}

\begin{twierdzenie}[K3: nieograniczona rodzina klas podstawowych]
\label{thm:tg-syn-k3}
Dla każdego $n\ge1$ niech $K_n=T(2,2n+1)$,
a $X_n=E(2)_{K_n}$ niech będzie wynikiem standardowej
chirurgii Fintushela--Sterna na włóknie eliptycznym K3.
Wówczas $X_n$ są homeomorficzne z K3, lecz nie są
dyfeomorficzne z K3; jeśli $m\ne n$,
to także $X_m$ i $X_n$ nie są dyfeomorficzne.
\end{twierdzenie}

\begin{proof}
Powierzchnia K3 ma $b_2^+=3$ i $\mathcal{SW}_{E(2)}=1$.
Twierdzenie Fintushela--Sterna daje więc
$\mathcal{SW}_{X_n}=\Delta_{K_n}(u)$,
$u=\exp(2[T_n])$. Dla węzła torusowego $K_n$,
\[
 \Delta_{K_n}(u)=\sum_{j=-n}^n(-1)^{n-j}u^j.
\]
Klasa włókna ma nieskończony rząd, więc wykładniki
$2j[T_n]$ są parami różne. Każdy ze współczynników
jest niezerowy, zatem $X_n$ ma dokładnie $2n+1$
klas podstawowych Seiberga--Wittena.
Liczba klas podstawowych jest niezmiennikiem
dyfeomorfizmu zachowującego orientację.
Dla różnych $n$ liczby te są różne, a dla K3
liczba ta wynosi jeden. Dyfeomorfizm odwracający
orientację jest wykluczony przez wspólną,
niezerową sygnaturę $-16$. Por. twierdzenie
\ref{thm:tg-ks-infinite}.
Homeomorfizm z K3 wynika z zachowania formy przecięcia,
prostej spójności i twierdzenia Freedmana, zgodnie
z poprzednią sekcją.
\end{proof}

\begin{uwaga}[Rachunek współczynników, nie tylko liczby klas]
Dla dwóch dowolnych węzłów liczba niezerowych współczynników
nie zawsze rozróżnia rozmaitości. Przy hipotezach twierdzenia \ref{thm:tg-ks-formula}
i ustalonej identyfikacji homologii, jeżeli
$\mathcal{SW}_X=\sum_\beta c_\beta\exp(\beta)$ oraz
$\Delta_K(u)=\sum_j a_ju^j$, to współczynnik przy
$\exp(\gamma)$ po chirurgii jest równy
\[
 \sum_{\beta+2j[T]=\gamma}c_\beta a_j.
\]
Różne pary $(\beta,j)$ mogą prowadzić do tej samej klasy
$\gamma$ i wtedy ich wkłady się sumują, niekiedy się znosząc.
Dlatego sama liczba wyrazów $\Delta_K$ w ogólnym przypadku
nie wyznacza liczby klas podstawowych $X_K$.
Jeżeli jednak $\mathcal{SW}_X=1$, to jedynym
niezerowym $c_\beta$ jest $c_0=1$ i ponieważ
$[T]$ ma nieskończony rząd, nie dochodzi do
sumowania wkładów różnych wykładników.
Odzyskujemy dokładnie stwierdzenie
\ref{prop:tg-ks-basic}.
\end{uwaga}

\section{Schemat hipotez i zależności między twierdzeniami}
\label{sec:tg-syn-assumptions}

Podajemy zestawienie ułatwiające rozpoznanie, jakie dane
należy sprawdzić, zanim zastosuje się określone twierdzenie.
Rzędy tabeli są niezależnymi implikacjami, nie automatycznym
łańcuchem równoważności.

\begin{center}
\begingroup
\footnotesize
\renewcommand{\arraystretch}{1.4}
\begin{tabular}{@{}p{.255\linewidth}p{.365\linewidth}p{.30\linewidth}@{}}
\hline
\textbf{Narzędzie} & \textbf{Kluczowe założenia} & \textbf{Wniosek}\\
\hline
{\raggedright Plumbing i dualność Poincar\'ego--Lefschetza\par}
 & {\raggedright zwarty orientowany blok, znana całkowita macierz przecięcia\par}
 & {\raggedright $H_1(\partial P)\cong\operatorname{coker}Q$ przy $H_1(P)=0$\par}\\
{\raggedright Wymierna kontrakcja\par}
 & {\raggedright konfiguracja $C_p$, zgodne sklejenie $B_p$, kontrola $\pi_1$\par}
 & {\raggedright zmiana $\chi,\sigma$; przy dalszych warunkach prosta spójność\par}\\
{\raggedright Freedman\par}
 & {\raggedright zamknięte, zorientowane, po prostu spójne topologiczne czterorozmaitości; izometria form i zgodność niezmienników Kirby'ego--Siebenmanna\par}
 & {\raggedright klasyfikacja topologiczna, nie gładka\par}\\
{\raggedright Fintushel--Stern\par}
 & {\raggedright zamknięta zorientowana gładka czterorozmaitość, $b_2^+>1$; regularne włókno typu \emph{cusp}, standardowe ramowanie, po prostu spójne dopełnienie\par}
 & {\raggedright formuła $\mathcal{SW}_{X_K}=\mathcal{SW}_X\Delta_K(u)$\par}\\
{\raggedright Taubes--Li--Liu\par}
 & {\raggedright zamknięta spójna symplektyczna czterorozmaitość; właściwa komora przy $b_2^+=1$\par}
 & {\raggedright niezmiennik kanonicznej struktury $\mathrm{Spin}^c$ równy $\pm1$ w komorze symplektycznej\par}\\
\hline
\end{tabular}
\endgroup
\end{center}

\begin{center}
\begin{tikzpicture}[>=Stealth,line cap=round,line join=round,
  every node/.style={align=center}]
  \node[draw=tgCoral,fill=tgCoral!10,rounded corners=4pt,
    text width=3.35cm,inner sep=6pt] (p0) at (-4.0,1.65)
    {\small $E(1)\#4\overline{\mathbb{CP}}^{\,2}$\\
     \footnotesize $C_7$, sześć sfer};
  \node[draw=tgCoral,fill=tgCoral!10,rounded corners=4pt,
    text width=3.35cm,inner sep=6pt] (p1) at (4.0,1.65)
    {\small $X_7$\\ \footnotesize $b_2^+=1$, $K\cdot\omega>0$};
  \node[draw=tgTeal,fill=tgTeal!10,rounded corners=4pt,
    text width=3.35cm,inner sep=6pt] (k0) at (-4.0,-1.65)
    {\small $E(2)$ (K3)\\ \footnotesize torus $T^2=0$};
  \node[draw=tgTeal,fill=tgTeal!10,rounded corners=4pt,
    text width=3.35cm,inner sep=6pt] (k1) at (4.0,-1.65)
    {\small $E(2)_K$\\ \footnotesize $b_2^+=3$, $\Delta_K(u)$};
  \draw[-{Stealth[length=2.7mm]},tgPurple,line width=1.7pt]
    (p0)--node[fill=white,font=\scriptsize,above=3pt]{wymierna kontrakcja}(p1);
  \draw[-{Stealth[length=2.7mm]},tgPurple,line width=1.7pt]
    (k0)--node[fill=white,font=\scriptsize,above=3pt]{chirurgia wzdłuż węzła}(k1);
  \node[text=tgNavy,font=\footnotesize] at (0,0)
    {ten sam język topologii, inne hipotezy teorii cechowania};
\end{tikzpicture}
\end{center}

\section{Granice przedstawionej syntezy}
\label{sec:tg-syn-limits}

Otrzymany obraz nie jest dowodem, że każdy węzeł albo każdy graf
plumbingowy daje egzotyczną czterorozmaitość. Potrzebne są
precyzyjne warunki istnienia zanurzeń i kontroli dopełnienia,
a także niezerowość właściwych niezmienników gładkich.

Z grafu plumbingowego można odczytać macierz przecięcia
\emph{określonego plumbingowania}, lecz nie wszystkie dane sklejenia
dowolnej rozmaitości. Z wielomianu Alexandera można obliczyć
pewne niezmienniki SW po chirurgii Fintushela--Sterna,
lecz dwa różne węzły mogą mieć ten sam wielomian.
Wymienione operacje są narzędziami konstrukcyjnymi,
a nie pełnymi klasyfikacjami w kategorii DIFF.

\begin{uwaga}[Jak korzystać z tego rozdziału]
Przy rozwiązywaniu konkretnego problemu należy przechodzić
w podanej kolejności: (1) dokładnie zdefiniować części i sklejenie;
(2) obliczyć homologię i grupę fundamentalną;
(3) wyznaczyć formę przecięcia i rozważyć typ homeomorfizmu;
(4) dopiero potem wybrać właściwe niezmienniki gładkie,
sprawdzając wszystkie ich hipotezy. Taka kolejność zapobiega
przeskokom od diagramu kombinatorycznego do wniosku o strukturze
gładkiej bez dowodu.
\end{uwaga}


\begin{uwaga}[Zakres wniosków klasyfikacyjnych]
Zgodność $e$, $\sigma$ i liczb Betti'ego nie klasyfikuje
czterorozmaitości. W przypadku rozmaitości po prostu
spójnych trzeba kontrolować całą całkowitą formę
przecięcia i odpowiednią hipotezę twierdzenia Freedmana.
Jeżeli rozmaitości mają niezerową grupę fundamentalną,
potrzebne są dalsze dane i twierdzenia klasyfikacyjne.
Niezmienniki Seiberga--Wittena badają strukturę gładką,
ale także nie stanowią pełnego układu niezmienników
klasyfikujących wszystkie czterorozmaitości.
\end{uwaga}

